% generated by gen_appA_r5.py from the Lean sources of the release (commit 4874cc0) and the census of commit 4874cc0; edit the generator, not this file.
\section{The Lean statements}\label{app:lean}
\sloppy
The statements below are copied from the Lean sources named before each block; only line breaks inside long lines are added by the
typesetting. They are quoted from the release of the repository \repourl\ (tag \repotag; Lean sources of commit \relcommit, \S\ref{ssec:formal}): the
family library from its library \texttt{ZetaShell}, the library of~\cite{DSouza} from its library \texttt{ZetaQ}, and the comparator files of
\S\ref{app:lean-release} from its folder \texttt{comparator/}. Each sha256 printed is that of the file as committed. Doc-strings are quoted as they
stand and are not part of the paper's claims: they are internal notes written during the formalisation, they refer to the working draft by its
labels and line numbers and to sources by internal abbreviations, they use the word ``level'' for the trust levels of the formalisation plan
(level A means kernel-checked in the sense of \S\ref{ssec:formal}, level C a displayed hypothesis), and some of them are not current. Only the Lean
code is claimed. Each status is that of the census of the release, taken on \censusref\ (\S\ref{app:lean-census}).

\subsection{Counting functions}
From \texttt{ZetaShell/Challenge.lean} (lines 31--90; sha256 \texttt{5bb7abc34ad3b557...}). This module imports only Mathlib. In these statements the height exponent is \texttt{r}
(our $r_0$).
\begin{Verbatim}[fontsize=\footnotesize,breaklines,frame=single,framesep=2mm]
open DirichletCharacter in
/-- `Σ*_{χ mod q}`: the primitive Dirichlet characters mod `q` (character-identical to `ZetaQ.primitiveChars`). -/
def primChars (q : ℕ) : Finset (DirichletCharacter ℂ q) :=
  letI := Classical.decPred (fun χ : DirichletCharacter ℂ q => χ.IsPrimitive)
  Finset.univ.filter (fun χ => χ.IsPrimitive)

section zeros
variable {q : ℕ} [NeZero q] (χ : DirichletCharacter ℂ q)

/-- `ρ` is a nontrivial zero of `L(s,χ)` (same as `Zeta23.ThmE.IsNontrivialZeroL`). -/
def IsNtZero (ρ : ℂ) : Prop := χ.LFunction ρ = 0 ∧ 0 < ρ.re ∧ ρ.re < 1

/-- `m_ρ`: order of vanishing of `L(·,χ)` at `ρ` (same as `Zeta23.ThmE.zeroMultL`). -/
def zmult (ρ : ℂ) : ℕ := (analyticOrderAt χ.LFunction ρ).toNat

/-- nontrivial zeros with `T₁ < γ ≤ T₂`. -/
def zerosWin (T₁ T₂ : ℝ) : Set ℂ := {ρ | IsNtZero χ ρ ∧ T₁ < ρ.im ∧ ρ.im ≤ T₂}

/-- `N_χ(T₁,T₂)`, with multiplicity. -/
def Nwin (T₁ T₂ : ℝ) : ℕ := ∑ᶠ ρ ∈ zerosWin χ T₁ T₂, zmult χ ρ

/-- `N^s_{0,χ}(T₁,T₂)`: simple zeros on the critical line. -/
def N0sWin (T₁ T₂ : ℝ) : ℕ :=
  (zerosWin χ T₁ T₂ ∩ {ρ | ρ.re = 1 / 2} ∩ {ρ | zmult χ ρ = 1}).ncard

/-- zeros in Jutila's / Bombieri's closed rectangle `σ ≤ β ≤ 1, |γ| ≤ T` (all zeros there are nontrivial
for `σ > 0`). -/
def zerosRect (σ T : ℝ) : Set ℂ := {ρ | χ.LFunction ρ = 0 ∧ σ ≤ ρ.re ∧ ρ.re ≤ 1 ∧ |ρ.im| ≤ T}

/-- `N(σ,T,χ)`, with multiplicity. -/
def NrectC (σ T : ℝ) : ℕ := ∑ᶠ ρ ∈ zerosRect χ σ T, zmult χ ρ

end zeros

/-- `N_χ(T₁,T₂)` as a real, with the (never used) guard `q = 0` (as `ZetaQ.NcountQ`). -/
def Nchi (q : ℕ) (χ : DirichletCharacter ℂ q) (T₁ T₂ : ℝ) : ℝ :=
  if h : q = 0 then 0 else haveI : NeZero q := ⟨h⟩; ((Nwin χ T₁ T₂ : ℕ) : ℝ)

/-- `N^s_{0,χ}(T₁,T₂)` as a real (as `ZetaQ.N0sQ`). -/
def N0schi (q : ℕ) (χ : DirichletCharacter ℂ q) (T₁ T₂ : ℝ) : ℝ :=
  if h : q = 0 then 0 else haveI : NeZero q := ⟨h⟩; ((N0sWin χ T₁ T₂ : ℕ) : ℝ)

/-- `N(σ,T,χ)` as a real. -/
def Nrect (q : ℕ) (χ : DirichletCharacter ℂ q) (σ T : ℝ) : ℝ :=
  if h : q = 0 then 0 else haveI : NeZero q := ⟨h⟩; ((NrectC χ σ T : ℕ) : ℝ)

/-- The moduli of the family `𝔉_Q` of sec_shell.tex l.37: `1 < q ≤ Q` (as `ZetaQ.Family.moduli Family.qle`). -/
def modQle (Q : ℕ) : Finset ℕ := Finset.Icc 2 Q

/-- The dyadic family's moduli `Q/2 < q ≤ Q` (as `ZetaQ.Family.moduli Family.dyadic`). -/
def modDyadic (Q : ℕ) : Finset ℕ := Finset.Ioc (Q / 2) Q

/-- `𝒩 = Σ_{q∈M} Σ*_χ N_χ(T₁,T₂)`. -/
def famN (M : Finset ℕ) (T₁ T₂ : ℝ) : ℝ := ∑ q ∈ M, ∑ χ ∈ primChars q, Nchi q χ T₁ T₂

/-- `Σ_{q∈M} Σ*_χ N^s_{0,χ}(T₁,T₂)`. -/
def famN0s (M : Finset ℕ) (T₁ T₂ : ℝ) : ℝ := ∑ q ∈ M, ∑ χ ∈ primChars q, N0schi q χ T₁ T₂

/-- The window height of the paper's Theorem 1: `T = (log Q)^{r+ε}` (as `ZetaQ.Twin`). -/
def twin (Q r ε : ℝ) : ℝ := Real.log Q ^ (r + ε)
\end{Verbatim}
These counts are those of the formal tree of~\cite{DSouza} by definition: from \texttt{ZetaShell/Frame/B1\_CountBridge.lean} (lines 10--22), five theorems proved by \texttt{rfl}
\tagL:
\begin{Verbatim}[fontsize=\footnotesize,breaklines,frame=single,framesep=2mm]
theorem famN_qle_eq (Qn : ℕ) (T₁ T₂ : ℝ) :
    famN (modQle Qn) T₁ T₂ = ZetaQ.NfamCount ZetaQ.Family.qle Qn T₁ T₂ := rfl

theorem famN0s_qle_eq (Qn : ℕ) (T₁ T₂ : ℝ) :
    famN0s (modQle Qn) T₁ T₂ = ZetaQ.N0sFamCount ZetaQ.Family.qle Qn T₁ T₂ := rfl

theorem famN_dyadic_eq (Qn : ℕ) (T₁ T₂ : ℝ) :
    famN (modDyadic Qn) T₁ T₂ = ZetaQ.NfamCount ZetaQ.Family.dyadic Qn T₁ T₂ := rfl

theorem famN0s_dyadic_eq (Qn : ℕ) (T₁ T₂ : ℝ) :
    famN0s (modDyadic Qn) T₁ T₂ = ZetaQ.N0sFamCount ZetaQ.Family.dyadic Qn T₁ T₂ := rfl

theorem twin_eq (Q r ε : ℝ) : twin Q r ε = ZetaQ.Twin Q r ε := rfl
\end{Verbatim}
and, from the library \texttt{ZetaQ}, \texttt{ZetaQ/Certificate.lean} (lines 570--575):
\begin{Verbatim}[fontsize=\footnotesize,breaklines,frame=single,framesep=2mm]
def NfamCount (F : Family) (Qn : ℕ) (T₁ T₂ : ℝ) : ℝ :=
  familySum F Qn (fun q χ => NcountQ q χ T₁ T₂)

/-- `Σ_{χ ∈ 𝔉_Q} N^s_{0,χ}(T₁,T₂)` — the left-hand side of Theorem 1 and of §3's display. -/
def N0sFamCount (F : Family) (Qn : ℕ) (T₁ T₂ : ℝ) : ℝ :=
  familySum F Qn (fun q χ => N0sQ q χ T₁ T₂)
\end{Verbatim}

\subsection{The zero-density input}\label{app:lean-density}
From \texttt{ZetaShell/Challenge.lean}, lines 94--124. \texttt{ZeroDensityInput} is (J) and (M) of \S\ref{ssec:shell-density}; it depends on no axiom. The
doc-string of \texttt{MontgomeryHybrid} is quoted in full.
The Prop itself states Bombieri's Th\'eor\`eme~20 with an existential power of the logarithm. Where the paper assumes (J) and (M), the Lean statement
has \texttt{hD : ZeroDensityInput} as a hypothesis; it is not an axiom, and \texttt{\#print axioms} does not list it.
\begin{Verbatim}[fontsize=\footnotesize,breaklines,frame=single,framesep=2mm]
/-- `N*(σ,T,Q) = Σ_{q ≤ Q} Σ*_{χ mod q} N(σ,T,χ)`, over PRIMITIVE characters, `q = 1` (i.e. `ζ`) INCLUDED
(sec_shell.tex l.742–744: "In every source below the count is over primitive characters, in the closed rectangle
`β ≥ σ, |γ| ≤ T`, and the character mod 1 (that is, ζ) is included"). -/
def Nstar (Q : ℕ) (σ T : ℝ) : ℝ := ∑ q ∈ Finset.Icc 1 Q, ∑ χ ∈ primChars q, Nrect q χ σ T

/-- **Jutila 1977, Theorem 1, (1.8)** [Math. Scand. 41 (1977), p. 46], as quoted in rem:shell-S53-inputs (a)
(sec_shell.tex l.791–796): "For `4/5 ≤ α ≤ 1`, `T ≥ 1`, we have … (1.8) `N*(α,T,Q) ≪_ε (Q²T)^{(2+ε)(1−α)}`",
"the constants implied by Vinogradov's symbol `≪` depend on `ε`"; no factor of `log`; no relation between `T` and `Q`.
Used as (LF) of (D1′) with `h_LF = 1`, `c₀ = 2`, `σ_LF = 4/5` (sec_shell.tex l.757).
Range of `Q`: `Q ≥ 1` (the sum over `q ≤ Q` is empty-free from `Q = 1`; TZ22 quote it for `Q ≥ 3`, which is
equivalent up to the constant since `N*` is monotone in `Q`). -/
def JutilaHybrid : Prop :=
  ∀ ε : ℝ, 0 < ε → ∃ C : ℝ, ∀ Q : ℕ, 1 ≤ Q → ∀ T : ℝ, 1 ≤ T → ∀ α : ℝ, 4 / 5 ≤ α → α ≤ 1 →
    Nstar Q α T ≤ C * ((Q : ℝ) ^ 2 * T) ^ ((2 + ε) * (1 - α))

/-- **Montgomery's hybrid estimate, in the form of Bombieri, Astérisque 18, §10, Théorème 20** (p. 77), as quoted in
rem:shell-S53-inputs (b) (sec_shell.tex l.802–807): for `T ≥ 2`, `1/2 ≤ α ≤ 1`,
"`Σ_{q≤Q} Σ*_{χ mod q} N(α,T;χ) ≪ (TQ²)^{3(1−α)/(2−α)} (log TQ)^A`", absolute implied constant.
The power `A` is EXISTENTIAL here.
(B) consumes only a finite `b₀`. So this Prop is implied by the printed statement for every admissible `A`, i.e.
it does NOT strengthen the source. -/
def MontgomeryHybrid : Prop :=
  ∃ A C : ℝ, ∀ Q : ℕ, 1 ≤ Q → ∀ T : ℝ, 2 ≤ T → ∀ α : ℝ, 1 / 2 ≤ α → α ≤ 1 →
    Nstar Q α T ≤ C * ((T * (Q : ℝ) ^ 2) ^ (3 * (1 - α) / (2 - α)) * Real.log (T * Q) ^ A)

/-- **`ZeroDensityInput`** (LEAN_PLAN §1.4): instantiation (D1′) of sec_shell.tex l.756–762 — "Montgomery bulk +
Jutila hybrid (all inputs primary, [V])" — which gives `A* = M = 5/2` and `α₀ = 5/3` (thm:shell-S53).
Both fields are classical unconditional theorems (Track D would prove them); here they are hypotheses. -/
structure ZeroDensityInput : Prop where
  jutila : JutilaHybrid
  montgomery : MontgomeryHybrid
\end{Verbatim}

\subsection{The certificate}
From \texttt{ZetaShell/Challenge.lean}, lines 139--229 (the profile \texttt{D53L75} and \texttt{CertD53} belong to the dyadic value of Remark~\ref{rem:shell-dyadic}, which is not
claimed):
\begin{Verbatim}[fontsize=\footnotesize,breaklines,frame=single,framesep=2mm]
/-- `ψ = v ⋆ v` read as the autocorrelation `ψ(α) = ∫ v(t) v(t−α) dt` (character-identical to
`ZetaQ.Payoff.psi`; for even `v` it is the convolution). -/
def psiS (v : ℝ → ℝ) (α : ℝ) : ℝ := ∫ t, v t * v (t - α)

/-- **The Shell kernel** (def:shell-kernel, sec_shell.tex l.56–59, in the normalisation of ssec:shell-cert
"Kernel and constants", l.1016–1018): `|α|` on `|α| ≤ 1`, the level `ℓ` on `1 < |α| ≤ α′`, the constant `C` on
`α′ < |α|` (only `|α| ≤ λ` is ever read, since `supp ψ ⊆ [−λ, λ]`). `ℓ = 1` for the full families. -/
def shellKernel (ℓ αp C : ℝ) (α : ℝ) : ℝ :=
  if |α| ≤ 1 then |α| else if |α| ≤ αp then ℓ else C

/-- **`B_{α′}(v) = ψ(0) + ∫ k_{α′}(α) ψ(α) dα`** (eq:shell-B with `k = k_{α′}`, def:shell-kernel). -/
def Bshell (ℓ αp C : ℝ) (v : ℝ → ℝ) : ℝ := psiS v 0 + ∫ α, shellKernel ℓ αp C α * psiS v α

/-- Admissible profiles (paper §11; field-identical to `ZetaQ.Payoff.Admissible`). -/
structure AdmissibleS (lam : ℝ) (v : ℝ → ℝ) : Prop where
  nonneg : ∀ t, 0 ≤ v t
  integrable : MeasureTheory.Integrable v
  mass : ∫ t, v t = 1
  supp : ∀ t, v t ≠ 0 → |t| ≤ lam / 2

/-- A certified profile of ssec:shell-cert "Method" (l.1025–1036): rationals `1 < λ < 2`, `α′`, the level `ℓ`, the
rational upper bound `C⁺` of the family constant, and `d = (d₀ = 1, d₁, …)` with `p(t) = Σ_i d_i (2t/λ)^{2i}`. -/
structure ShellProfile where
  lam : ℚ
  alphaP : ℚ
  level : ℚ
  Cplus : ℚ
  d : List ℚ

namespace ShellProfile
variable (S : ShellProfile)

/-- `p(t) = Σ_i d_i (2t/λ)^{2i}`. -/
def p (t : ℝ) : ℝ := ∑ i : Fin S.d.length, ((S.d.get i : ℚ) : ℝ) * (2 * t / (S.lam : ℝ)) ^ (2 * (i : ℕ))

/-- `∫_{−λ/2}^{λ/2} p²`. -/
def mass : ℝ := ∫ t in (-((S.lam : ℝ) / 2))..((S.lam : ℝ) / 2), S.p t ^ 2

/-- `v_p = p² / ∫p²` on `[−λ/2, λ/2]`, `0` outside. -/
def v (t : ℝ) : ℝ := if |t| ≤ (S.lam : ℝ) / 2 then S.p t ^ 2 / S.mass else 0

/-- **The Lean side conditions "L75"** (ssec:shell-cert "Lean side conditions", l.1080–1082: `p` decreasing,
`a ≥ 3/4` (`Valid.a_ge`), `b ≥ 1/2` (H5), `p(λ/2) ≥ 1/6`), with `0 < p ≤ 1` of "Method" (l.1033–1034);
`a = ⟨p²⟩ = λ⁻¹∫p²`, `b = ⟨p⁴⟩ = λ⁻¹∫p⁴` as in `check_profile.py`. -/
structure L75 : Prop where
  pos : ∀ t, |t| ≤ (S.lam : ℝ) / 2 → 0 < S.p t
  le_one : ∀ t, |t| ≤ (S.lam : ℝ) / 2 → S.p t ≤ 1
  anti : StrictAntiOn S.p (Set.Icc 0 ((S.lam : ℝ) / 2))
  a_ge : (3 : ℝ) / 4 ≤ S.mass / (S.lam : ℝ)
  b_ge : (1 : ℝ) / 2 ≤ (∫ t in (-((S.lam : ℝ) / 2))..((S.lam : ℝ) / 2), S.p t ^ 4) / (S.lam : ℝ)
  end_ge : (1 : ℝ) / 6 ≤ S.p ((S.lam : ℝ) / 2)

end ShellProfile

/-- **The certificate Prop (Track R)**: the profile is admissible at `λ`, meets L75, `1 < α′`, `1 < λ < 2`, and
`B_{α′}(v_p) ≤ 2 − P_c` at the kernel `(ℓ, α′, C⁺)`. -/
def ShellCert (S : ShellProfile) (Pc : ℚ) : Prop :=
  1 < (S.lam : ℝ) ∧ (S.lam : ℝ) < 2 ∧ 1 < (S.alphaP : ℝ) ∧ AdmissibleS S.lam S.v ∧ S.L75 ∧
    Bshell S.level S.alphaP S.Cplus S.v ≤ 2 - (Pc : ℝ)

/-- Profile **S53-L75** (ssec:shell-cert-53, l.1224; `round2/d2_1_numerics/profiles53.json`):
`α′ = 2497/1500 = 5/3 − 1/500`, `λ = 191/100`, `ℓ = 1`, `C⁺ = 541161617/10⁸ ≥ π⁴/18`. -/
def S53L75 : ShellProfile where
  lam := 191 / 100
  alphaP := 2497 / 1500
  level := 1
  Cplus := 541161617 / 100000000
  d := [1, -109245641 / 315967035, -484609675 / 903872491, 1126440115 / 499347366,
        -1244916364 / 475100341, 381673291 / 981251350, 37853428 / 698475165]

/-- Profile **D53-L75** (l.1230): `α′ = 2497/1500`, `λ = 93/50`, `ℓ = 1`, `C⁺ = 721548823/10⁸ ≥ 2π⁴/27`. -/
def D53L75 : ShellProfile where
  lam := 93 / 50
  alphaP := 2497 / 1500
  level := 1
  Cplus := 721548823 / 100000000
  d := [1, -390353034 / 798066593, 1309177511 / 918253499, -6068806879 / 896328244,
        13364613432 / 870778397, -12836034741 / 854019032, 4242100897 / 905838731]

/-- `0.9059137927` — the 10-digit truncation of `2 − B(v_{S53-L75})` (exact rational `0.90591379273786…`,
`p_cert_exact53.txt`). -/
def PcertS53 : ℚ := 9059137927 / 10000000000

/-- `0.9031776196` — the 10-digit truncation for D53-L75 (exact `0.90317761969178…`). -/
def PcertD53 : ℚ := 9031776196 / 10000000000

/-- The certificate of the headline (sharp family). -/
def CertS53 : Prop := ShellCert S53L75 PcertS53

/-- The certificate of the dyadic headline. -/
def CertD53 : Prop := ShellCert D53L75 PcertD53
\end{Verbatim}
From \texttt{ZetaShell/Cert/R4\_CertS53.lean} (sha256 \texttt{6353961802d5fc7d...}), the two statements (proofs omitted):
\begin{Verbatim}[fontsize=\footnotesize,breaklines,frame=single,framesep=2mm]
theorem certS53 : CertS53
theorem certD53 : CertD53
\end{Verbatim}
Status: both kernel-checked, no \texttt{sorry}, no \texttt{native\_decide}; in the census both depend on the three standard axioms only and rest
on no open declaration \tagL.

\subsection{The statement of record of Theorem~\ref{thm:shell}}\label{app:lean-record}
From \texttt{ZetaShell/Top/HeadlineReduced.lean} (sha256 \texttt{e4bb7a2ca9f41c53...}), lines 37--47, with its proof:
\begin{Verbatim}[fontsize=\footnotesize,breaklines,frame=single,framesep=2mm]
/-- **Reduced headline, family `1 < q ≤ Q`** (statement of record): `shell_S53_qle` without `hP` and `hcert`. -/
theorem shell_S53_qle_reduced (hD : ZeroDensityInput) (r ε θ : ℝ) (hr : 3 ≤ r) (hε : 0 < ε)
    (hθ : 0 < θ) (hθ' : θ < 503 / 1994) :
    ∃ Q₀ c : ℝ, 0 < c ∧ ∀ Qn : ℕ, Q₀ ≤ (Qn : ℝ) →
      (((PcertS53 : ℚ) : ℝ) - c * Real.log (Qn : ℝ) ^ (-θ))
          * famN (modQle Qn) (twin (Qn : ℝ) r ε) (2 * twin (Qn : ℝ) r ε)
        ≤ famN0s (modQle Qn) (twin (Qn : ℝ) r ε) (2 * twin (Qn : ℝ) r ε) := by
  obtain ⟨Q₀, c, hc, h⟩ := shell_S53_qle_tree_noP hD r ε θ hr hε hθ hθ'
  refine ⟨Q₀, c, hc, fun Qn hQn => ?_⟩
  rw [famN_qle_eq, famN0s_qle_eq, twin_eq]
  exact h Qn hQn
\end{Verbatim}
In the release this is a theorem \tagL: its row in the census table (tab-separated; the tabs are printed here as two spaces) lists no open leaf,
and \texttt{\#print axioms} gives the three standard axioms only:
\begin{Verbatim}[fontsize=\footnotesize,breaklines,frame=single,framesep=2mm]
ZetaShell.Top.HeadlineReduced  ZetaShell.shell_S53_qle_reduced  A  [Classical.choice, Quot.sound, propext]  0  []
'ZetaShell.shell_S53_qle_reduced' depends on axioms: [propext, Classical.choice, Quot.sound]
\end{Verbatim}
Its one displayed hypothesis, \texttt{hD : ZeroDensityInput}, is (J) and (M) (\S\ref{app:lean-density}). The proof applies
\texttt{shell\_S53\_qle\_tree\_noP} of \texttt{ZetaShell/Top/T1\_HeadlineNoP.lean}, the headline obtained from the frame
\texttt{shell\_frame\_qle\_of\_K\_noP}, which does not use \texttt{PNTErrorTerm}, and rewrites the counts by the theorems of the first subsection that are
proved by \texttt{rfl}.
The same module proves, as \texttt{ZetaShell.Top.shell\_S53\_qle\_of\_reduced} (lines 91--93, no \texttt{sorry} of its own), that the
statement of record implies the earlier form \texttt{ZetaShell.shell\_S53\_qle} of \texttt{ZetaShell/Challenge.lean}, which has the two further
hypotheses \texttt{hP : PNTErrorTerm} and \texttt{hcert : CertS53}. The converse is proved given \texttt{PNTErrorTerm}; the tree does not prove
\texttt{PNTErrorTerm}, and by the module's header the argument uses only the medium form of the prime number theorem, which is a theorem of the tree:
\begin{Verbatim}[fontsize=\footnotesize,breaklines,frame=single,framesep=2mm]
theorem shell_S53_qle_of_reduced :
    type_of% @shell_S53_qle_reduced → type_of% @shell_S53_qle :=
  fun h hD _ _ => h hD
\end{Verbatim}
The same files also state a dyadic version (\texttt{shell\_S53\_dyadic\_reduced}) and two bounded-height versions (\texttt{shell\_S53\_bounded\_reduced},
\texttt{shell\_S53\_bounded\_dyadic\_reduced}); they are open, and none of them is claimed in this paper.

\subsection{The statement of Theorem~\ref{thm:one-prime}}
From \texttt{ZetaShell/Defs/LK\_Defs.lean} (sha256 \texttt{d05d5949ca23b6f3...}), lines 104--105:
\begin{Verbatim}[fontsize=\footnotesize,breaklines,frame=single,framesep=2mm]
/-- the constant of the design-profile route (no new design data): `0.7235`. -/
def PcertKD : ℚ := 7235 / 10000
\end{Verbatim}
and from \texttt{ZetaShell/LemmaK/LK\_KT\_Headline.lean} (sha256 \texttt{3cf2934de48e80cc...}), lines 60--65 (proof omitted):
\begin{Verbatim}[fontsize=\footnotesize,breaklines,frame=single,framesep=2mm]
/-- **THEOREM 1′ at ZetaQ's own design profile, `P = 0.7235`.** -/
theorem theorem_one_prime_design (r ε : ℝ) (hr : 3 ≤ r) (hε : 0 < ε) :
    ∃ Q₀ c : ℝ, 0 < c ∧ ∀ Qn : ℕ, Q₀ ≤ (Qn : ℝ) →
      (((PcertKD : ℚ) : ℝ) - c * Real.log (Real.log (Qn : ℝ)) / Real.log (Qn : ℝ))
          * ZetaQ.NfamCount ZetaQ.Family.qle Qn (ZetaQ.Twin (Qn : ℝ) r ε) (2 * ZetaQ.Twin (Qn : ℝ) r ε)
        ≤ ZetaQ.N0sFamCount ZetaQ.Family.qle Qn (ZetaQ.Twin (Qn : ℝ) r ε) (2 * ZetaQ.Twin (Qn : ℝ) r ε)
\end{Verbatim}
Status: a Lean theorem with the three standard axioms and no hypothesis, in the library \texttt{ZetaShell} of the release \tagL. The census
gives
\begin{Verbatim}[fontsize=\footnotesize,breaklines,frame=single,framesep=2mm]
SORRY-LEAVES ZetaShell.LemmaK.theorem_one_prime_design (0): []
'ZetaShell.LemmaK.theorem_one_prime_design' depends on axioms: [propext, Classical.choice, Quot.sound]
\end{Verbatim}
and so does the file \texttt{comparator/PrintAxioms/FollowUpShell.lean} (\S\ref{app:lean-release}). The companion statement at $0.7237$
(\texttt{theorem\_one\_prime}) still rests on one open node (\texttt{killed\_frame\_lamStar8}). The statement has the shape of the theorem of~\cite{DSouza} in the library
\texttt{ZetaQ}, which is kernel-checked with no hypothesis; from \texttt{ZetaQ/Margin.lean} (lines 1215--1220; proof omitted):
\begin{Verbatim}[fontsize=\footnotesize,breaklines,frame=single,framesep=2mm]
theorem theorem_one_generic_proved' (r ε : ℝ) (hr : 3 ≤ r) (hε : 0 < ε) :
    ∃ Q₀ c : ℝ, 0 < c ∧ ∀ Qn : ℕ, Q₀ ≤ (Qn : ℝ) →
      (((Payoff.Pcert_qQ_smooth : ℚ) : ℝ)
          - c * Real.log (Real.log (Qn : ℝ)) / Real.log (Qn : ℝ))
          * NfamCount Family.qle Qn (Twin (Qn : ℝ) r ε) (2 * Twin (Qn : ℝ) r ε)
        ≤ N0sFamCount Family.qle Qn (Twin (Qn : ℝ) r ε) (2 * Twin (Qn : ℝ) r ε)
\end{Verbatim}
Here \texttt{Payoff.Pcert\_qQ\_smooth} is the constant $0.7212$ of~\cite{DSouza}.

\subsection{The census, and the parts of Theorem~\ref{thm:shell}}\label{app:lean-census}
The census is one Lean job that prints, for each public theorem of \texttt{ZetaShell}, the axioms it depends on and the open declarations
(declarations with a \texttt{sorry} of their own) that its proof term reaches. It was run on \censusref\ (output:
\texttt{audit/followup/census\_2026-10-03\_1557.out} and \texttt{audit/followup/census\_2026-10-03\_1557.tsv} in the repository; the job also covers \texttt{ZetaS}): of the 1,744 public theorems, 1,708 depend on the three standard axioms only and 36 on \texttt{sorryAx}: 29 open
statements, each with its own \texttt{sorry} (20 in \texttt{ZetaShell/Skeleton/}, and the 9 headline statements of \texttt{ZetaShell/Challenge.lean},
\texttt{ZetaShell/Interfaces.lean} and \texttt{ZetaShell/Top/HeadlineReduced.lean}), and 7 theorems derived from them. No other axiom occurs.
For the parts of Theorem~\ref{thm:shell} named in \S\ref{ssec:formal} the census prints (records of the open leaves first, then the axioms; verbatim):
\begin{Verbatim}[fontsize=\footnotesize,breaklines,frame=single,framesep=2mm]
SORRY-LEAVES ZetaShell.shell_S53_qle_reduced (0): []
SORRY-LEAVES ZetaShell.Design.shell_frame_qle_of_K_noP (0): []
SORRY-LEAVES ZetaShell.Design.shell_frame_qle_of_K (0): []
SORRY-LEAVES ZetaShell.ShellK.thmK (0): []
SORRY-LEAVES ZetaShell.ShellK.K2_shellZone (0): []
SORRY-LEAVES ZetaShell.ShellK.KT_transfer (0): []
SORRY-LEAVES ZetaShell.ShellS.AS_pointwise_corr (0): []
SORRY-LEAVES ZetaShell.ShellK.F1c.K3b_integrated (0): []
SORRY-LEAVES ZetaShell.ShellS.shell_S (0): []
SORRY-LEAVES ZetaShell.ShellS.S1_ring_to_zeros (0): []
SORRY-LEAVES ZetaShell.ShellS.S2_zero_side (0): []
SORRY-LEAVES ZetaShell.ShellS.Z6a_zero_sums (0): []
SORRY-LEAVES ZetaShell.ShellS.Z6d_counts_corr (0): []
SORRY-LEAVES ZetaShell.TrackF.ring_le_primitive_corr (0): []
SORRY-LEAVES ZetaShell.TrackF.lemma2a (0): []
SORRY-LEAVES ZetaShell.TrackF.lemmaP (0): []
SORRY-LEAVES ZetaShell.PropZ.propZ_W (0): []
SORRY-LEAVES ZetaShell.ShellK.F1c_chain (0): []
SORRY-LEAVES ZetaShell.TrackF.signed_gallagher (0): []
SORRY-LEAVES ZetaShell.ShellS.S3_family_errors (0): []
SORRY-LEAVES ZetaShell.certS53 (0): []
'ZetaShell.shell_S53_qle_reduced' depends on axioms: [propext, Classical.choice, Quot.sound]
'ZetaShell.Design.shell_frame_qle_of_K_noP' depends on axioms: [propext, Classical.choice, Quot.sound]
'ZetaShell.Design.shell_frame_qle_of_K' depends on axioms: [propext, Classical.choice, Quot.sound]
'ZetaShell.ShellK.thmK' depends on axioms: [propext, Classical.choice, Quot.sound]
'ZetaShell.ShellK.K2_shellZone' depends on axioms: [propext, Classical.choice, Quot.sound]
'ZetaShell.ShellK.KT_transfer' depends on axioms: [propext, Classical.choice, Quot.sound]
'ZetaShell.ShellS.AS_pointwise_corr' depends on axioms: [propext, Classical.choice, Quot.sound]
'ZetaShell.ShellK.F1c.K3b_integrated' depends on axioms: [propext, Classical.choice, Quot.sound]
'ZetaShell.ShellS.shell_S' depends on axioms: [propext, Classical.choice, Quot.sound]
'ZetaShell.ShellS.S1_ring_to_zeros' depends on axioms: [propext, Classical.choice, Quot.sound]
'ZetaShell.ShellS.S2_zero_side' depends on axioms: [propext, Classical.choice, Quot.sound]
'ZetaShell.ShellS.Z6a_zero_sums' depends on axioms: [propext, Classical.choice, Quot.sound]
'ZetaShell.ShellS.Z6d_counts_corr' depends on axioms: [propext, Classical.choice, Quot.sound]
'ZetaShell.TrackF.ring_le_primitive_corr' depends on axioms: [propext, Classical.choice, Quot.sound]
'ZetaShell.TrackF.lemma2a' depends on axioms: [propext, Classical.choice, Quot.sound]
'ZetaShell.TrackF.lemmaP' depends on axioms: [propext, Classical.choice, Quot.sound]
'ZetaShell.PropZ.propZ_W' depends on axioms: [propext, Classical.choice, Quot.sound]
'ZetaShell.ShellK.F1c_chain' depends on axioms: [propext, Classical.choice, Quot.sound]
'ZetaShell.TrackF.signed_gallagher' depends on axioms: [propext, Classical.choice, Quot.sound]
'ZetaShell.ShellS.S3_family_errors' depends on axioms: [propext, Classical.choice, Quot.sound]
'ZetaShell.certS53' depends on axioms: [propext, Classical.choice, Quot.sound]
\end{Verbatim}
So in the release each of these depends on the three standard axioms only and its proof reaches no open declaration: the statement of record,
the frame (also without \texttt{PNTErrorTerm}), Theorem~\ref{thm:shell-K} with its Shell-zone part, the transfer, the pointwise assembly in its corrected
form and the integrated killed-zone part, Theorem~\ref{thm:shell-S} with its steps S1 and S2 and Lemmas~\ref{lem:shell-6a} and~\ref{lem:shell-6d} (the
latter corrected), the corrected Lemma~\ref{lem:shell-star}, Lemma~\ref{lem:shell-2}(a), Lemma~\ref{lem:shell-P}, Proposition~\ref{prop:shell-Z}, the
Frobenius-row chain, Lemma~\ref{lem:shell-3}, the family errors of Theorem~\ref{thm:shell-S} and the certificate are theorems. The table below gives
the status of each named node.

\subsection{Status of the named nodes}
The label of each row is computed from the census of \S\ref{app:lean-census}.
For \tagLm\ the open leaves are listed.
\begin{center}\footnotesize
\begin{tabular}{@{}>{\raggedright\arraybackslash}p{5.0cm}>{\raggedright\arraybackslash}p{4.6cm}>{\raggedright\arraybackslash}p{5.0cm}@{}}
\toprule
declaration & paper & status\\
\midrule
\texttt{certS53} & \S\ref{sec:cert} & \tagL\\
\texttt{certD53} & Remark~\ref{rem:shell-dyadic} (not claimed) & \tagL\\
\texttt{psi\_sub\_id\_isLittleO\_log\_rpow} & \eqref{eq:PNT} & \tagL\\
\texttt{signed\_farey\_identity} & Lemma~\ref{lem:shell-1} & \tagL\\
\texttt{replace\_cost}, \texttt{small\_prime\_norm}, \texttt{far\_norm} & Lemma~\ref{lem:shell-1prime} & \tagL\\
\texttt{hole\_edge\_mass}, \texttt{shell\_rectangles} & Lemmas~\ref{lem:shell-F1}, \ref{lem:shell-6b} & \tagL\\
\texttt{lemma2b}, \texttt{lemma2c\_hole}, \texttt{lemma2c\_spike} & Lemma~\ref{lem:shell-2}(b),(c) & \tagL\\
\texttt{lemma2a} & Lemma~\ref{lem:shell-2}(a) & \tagL\\
\texttt{lemmaP} & Lemma~\ref{lem:shell-P} & \tagL\\
\texttt{signed\_gallagher} & Lemma~\ref{lem:shell-3} & \tagL\\
\texttt{ring\_le\_primitive\_corr} & Lemma~\ref{lem:shell-star} (corrected statement, Appendix~\ref{app:corr}) & \tagL\\
\texttt{plancherel\_majorant} & Lemma~\ref{lem:shell-5a} & \tagL\\
\texttt{smooth\_explicit\_formula\_weil} & Lemma~\ref{lem:shell-5b} & \tagL\\
\texttt{family\_E\_sum}, \texttt{family\_EW\_sum} & Lemma~\ref{lem:shell-5c}(2),(2$'$) & \tagL\\
\texttt{propZ\_W} & Proposition~\ref{prop:shell-Z} & \tagL\\
\texttt{Z6c\_blocks} & Lemma~\ref{lem:shell-6c} & \tagL\\
\texttt{Z6a\_zero\_sums} & Lemma~\ref{lem:shell-6a} & \tagL\\
\texttt{Z6d\_counts\_corr} & Lemma~\ref{lem:shell-6d} (near bound under $5X_0\le s_0$, Appendix~\ref{app:corr}) & \tagL\\
\texttt{Z6d\_counts} & first form of Lemma~\ref{lem:shell-6d} (not used; Appendix~\ref{app:corr}) & \tagLs, open\\
\texttt{S1\_ring\_to\_zeros} & Theorem~\ref{thm:shell-S}, step S1 & \tagL\\
\texttt{S2\_zero\_side} & Theorem~\ref{thm:shell-S}, step S2 & \tagL\\
\texttt{S3\_family\_errors}, \texttt{S4\_elementary} & Theorem~\ref{thm:shell-S}, steps S3, S4 & \tagL\\
\texttt{shell\_S} & Theorem~\ref{thm:shell-S} & \tagL\\
\texttt{AS\_pointwise\_corr} & \eqref{eq:shell-assembly} (with $\alpha''<\lambda$, Appendix~\ref{app:corr}) & \tagL\\
\texttt{AS\_pointwise} & first form of \eqref{eq:shell-assembly} (not used; Appendix~\ref{app:corr}) & \tagLs, open\\
\texttt{K1\_strip}, \texttt{K\_int} & Theorem~\ref{thm:shell-K}: strip, integrability & \tagL\\
\texttt{killed\_pointwise\_shell}, \texttt{majorant\_excess} & Theorem~\ref{thm:shell-K}: killed zone pointwise (K3a), majorant & \tagL\\
\texttt{K2\_shellZone}, \texttt{K3b\_integrated}, \texttt{KT\_transfer} & Theorem~\ref{thm:shell-K}: Shell zone, killed zone integrated (K3b), transfer & \tagL\\
\texttt{thmK} & Theorem~\ref{thm:shell-K} & \tagL\\
\texttt{F1c\_chain} & Frobenius-row chain & \tagL\\
\texttt{zoneLipschitzData\_S53} & zone slope, \S\ref{sec:design}(viii) (not used) & \tagLs, open\\
\texttt{shell\_frame\_qle\_of\_nodes} & frame of Theorem~\ref{thm:shell}, Frobenius row as one node & \tagLm, one leaf: \texttt{frob\_row\_shell}\\
\texttt{shell\_frame\_qle\_of\_K} & frame of Theorem~\ref{thm:shell}, Frobenius row from Theorem~\ref{thm:shell-K} & three standard axioms, no open leaf, with the further hypotheses \texttt{PNTErrorTerm} and \texttt{CertS53}\\
\texttt{shell\_frame\_qle\_of\_K\_noP} & the same frame without \texttt{PNTErrorTerm} & \tagL\\
\texttt{gauss\_sum\_ineq}, \texttt{farey\_hole\_cores} & Lemmas~\ref{lem:K-farey}, \ref{lem:K-Z} & \tagL\\
\texttt{gallagher\_with\_zero\_set\_corr} & Lemma~\ref{lem:K-G} & \tagL\\
\texttt{principal\_gauss\_hole}, \texttt{small\_primes\_l2} & Lemmas~\ref{lem:K-R}, \ref{lem:K-S} & \tagL\\
\texttt{hole\_lower}, \texttt{farey\_hole\_bound} & \eqref{eq:K-hole} & \tagL\\
\texttt{killed\_pointwise} & Lemma~\ref{lem:K}(i) & \tagL\\
\texttt{principal\_arc\_mass} & Lemma~\ref{lem:K-P} & \tagL\\
\texttt{certKD\_core}, \texttt{k\_assembly} & \S\ref{ssec:K-constants}, \S\ref{ssec:thm1prime} & \tagL\\
\texttt{theorem\_one\_prime\_design} & Theorem~\ref{thm:one-prime} & \tagL\\
\texttt{shell\_S53\_qle\_reduced} & Theorem~\ref{thm:shell} & \tagL\ (\S\ref{app:lean-record})\\
\bottomrule
\end{tabular}
\end{center}

\subsection{The release, and how to check it}\label{app:lean-release}
The formalisation is released in the repository \repourl, at tag \repotag; its Lean sources are those of commit \relcommit\ of the
development history, as \texttt{RELEASE\_NOTES.md} of the repository records (the repository is a snapshot of the release). The family library of this paper is its library
\texttt{ZetaShell} (327 files, 51,318 lines; the groups of modules are described in \texttt{ZetaShell/README.md}); the repository also holds the Lean artifact \texttt{Zeta23} of~\cite{AF}, the library
\texttt{ZetaQ} of~\cite{DSouza} and the libraries of the companion paper on~$\zeta$. Commands are run from the root of the repository; the toolchain
is pinned in \texttt{lean-toolchain}, and \texttt{lake exe cache get} fetches the prebuilt Mathlib. Then
\begin{Verbatim}[fontsize=\small,xleftmargin=1em]
lake build ZetaShell
lake env lean comparator/PrintAxioms/FollowUpShell.lean
\end{Verbatim}
The last command prints the axioms of \texttt{ZetaShell.LemmaK.theorem\_one\_prime\_design} and \texttt{ZetaShell.certS53}, and must print the two lines
\begin{Verbatim}[fontsize=\footnotesize,xleftmargin=1em]
'ZetaShell.LemmaK.theorem_one_prime_design' depends on axioms: [propext, Classical.choice, Quot.sound]
'ZetaShell.certS53' depends on axioms: [propext, Classical.choice, Quot.sound]
\end{Verbatim}
The build warns that a declaration uses \texttt{sorry} for each of the open statements of \S\ref{app:lean-census}; whether a theorem depends on an
open statement is decided by \texttt{\#print axioms}, not by these warnings. The build of the release commit had
no error and printed these two lines; its log and outputs, and those of earlier builds and censuses, are in \texttt{audit/followup/} of the
repository, and \texttt{RELEASE\_NOTES.md} names the commit each was taken on.

The two theorems of this paper are also stated in the comparator package \texttt{FollowUpFamily}: the trusted file \texttt{comparator/ChallengeDeps/FollowUpFamily.lean} (sha256
\texttt{07f36f808a8c6f35...}) imports Mathlib alone and defines the counts, the family, $T$ and \texttt{ZeroDensityInput} from their sources, and the trusted file
\texttt{comparator/Challenge/FollowUpFamily.lean} (sha256 \texttt{5cc2d104721ecafe...}) states the two theorems, lines 25--47 (the \texttt{sorry} is the challenge side):
\begin{Verbatim}[fontsize=\footnotesize,breaklines,frame=single,framesep=2mm]
/-- **Theorem thm:one-prime (killed kernel)** of the family paper: for `T = (log Q)^{r+ε}`, `r ≥ 3`, `ε > 0`, there are
`Q₀` and `c > 0` such that for every integer `Q ≥ Q₀`
`(0.7235 − c·log log Q/log Q) · Σ_{1<q≤Q} Σ*_χ N_χ(T,2T) ≤ Σ_{1<q≤Q} Σ*_χ N^s_{0,χ}(T,2T)`.
No hypothesis. (In the library: `ZetaShell.LemmaK.theorem_one_prime_design`.) -/
theorem family_simple_on_line_killed (r ε : ℝ) (hr : 3 ≤ r) (hε : 0 < ε) :
    ∃ Q₀ c : ℝ, 0 < c ∧ ∀ Qn : ℕ, Q₀ ≤ (Qn : ℝ) →
      ((7235 : ℝ) / 10 ^ 4 - c * Real.log (Real.log (Qn : ℝ)) / Real.log (Qn : ℝ))
          * famN (modQle Qn) (twin (Qn : ℝ) r ε) (2 * twin (Qn : ℝ) r ε)
        ≤ famN0s (modQle Qn) (twin (Qn : ℝ) r ε) (2 * twin (Qn : ℝ) r ε) := by
  sorry

/-- **Theorem thm:shell (the Shell kernel)** of the family paper: assume the zero-density estimates (J) and (M)
(`ZeroDensityInput`). For `T = (log Q)^{r+ε}`, `r ≥ 3`, `ε > 0`, and `0 < θ < 503/1994`, there are `Q₀` and `c > 0`
such that for every integer `Q ≥ Q₀`
`(0.9059137927 − c·(log Q)^{−θ}) · Σ_{1<q≤Q} Σ*_χ N_χ(T,2T) ≤ Σ_{1<q≤Q} Σ*_χ N^s_{0,χ}(T,2T)`.
(In the library: `ZetaShell.shell_S53_qle_reduced`.) -/
theorem family_simple_on_line_shell (hD : ZeroDensityInput) (r ε θ : ℝ) (hr : 3 ≤ r) (hε : 0 < ε)
    (hθ : 0 < θ) (hθ' : θ < 503 / 1994) :
    ∃ Q₀ c : ℝ, 0 < c ∧ ∀ Qn : ℕ, Q₀ ≤ (Qn : ℝ) →
      ((9059137927 : ℝ) / 10 ^ 10 - c * Real.log (Qn : ℝ) ^ (-θ))
          * famN (modQle Qn) (twin (Qn : ℝ) r ε) (2 * twin (Qn : ℝ) r ε)
        ≤ famN0s (modQle Qn) (twin (Qn : ℝ) r ε) (2 * twin (Qn : ℝ) r ε) := by
  sorry
\end{Verbatim}
They are proved in \texttt{comparator/Solution/FollowUpFamily.lean} from \texttt{ZetaShell.LemmaK.theorem\_one\_prime\_design} and
\texttt{ZetaShell.shell\_S53\_qle\_reduced}. To check them, run, after \texttt{lake exe cache get},
\begin{Verbatim}[fontsize=\small,xleftmargin=1em]
lake build Challenge.FollowUpFamily Solution.FollowUpFamily
lake env lean comparator/PrintAxioms/FollowUpFamily.lean
\end{Verbatim}
The last command must print
\begin{Verbatim}[fontsize=\footnotesize,xleftmargin=1em]
'family_simple_on_line_killed' depends on axioms: [propext, Classical.choice, Quot.sound]
'family_simple_on_line_shell' depends on axioms: [propext, Classical.choice, Quot.sound]
\end{Verbatim}
as it did in the build of the release commit (output in
\texttt{audit/followup/}); \texttt{\#print axioms ZetaShell.shell\_S53\_qle\_reduced} in a
file that imports \texttt{ZetaShell} prints the line quoted in \S\ref{app:lean-record}. The census files of \S\ref{app:lean-census} are in \texttt{audit/followup/} as well.
